\section{结果文件与复现说明}

核心脚本及输出关系见\cref{tab:appendix-files}。所有脚本从项目根目录解析路径，计算结果分别写入 \texttt{results/q1} 至 \texttt{results/q5}，论文图写入 \texttt{paper/figures}。

\begin{table}[!htbp]
 \centering
 \caption{主要程序和结果文件}
 \label{tab:appendix-files}
 \begin{tabular}{ll}
  \toprule[1.5pt]
  文件 & 作用 \\
  \midrule[1pt]
  \texttt{q1\_smoke\_duration.py} & 问题一给定策略计算与灵敏度分析 \\
  \texttt{q2\_single\_smoke\_optimization.py} & 问题二单弹多盆地优化 \\
  \texttt{q3\_three\_smoke\_optimization.py} & 问题三共享航迹三弹优化 \\
  \texttt{q4\_q5\_joint\_optimization.py} & 问题四、五协同求解与终算 \\
  \texttt{result1.xlsx--result3.xlsx} & 问题三至五官方格式结果表 \\
  \bottomrule[1.5pt]
 \end{tabular}
\end{table}

在已安装 NumPy、SciPy 和 Matplotlib 的 Python 环境中，可在项目根目录执行相应脚本。最终论文使用 XeLaTeX 连续编译两遍，以更新交叉引用和页码。Excel 文件由同一份 JSON 结果自动填入官方模板，因此不需要人工复制坐标。
\endinput
\section{蒙特卡洛算法--matlab 代码}

\begin{lstlisting}[language=matlab]
% 蒙特卡洛估算圆周率
N = 1e6; % 随机点数量
x = rand(N,1);
y = rand(N,1);
inside = (x.^2 + y.^2) <= 1;
pi_est = 4 * sum(inside) / N;
fprintf('估算的值: %.6f\n', pi_est);
\end{lstlisting}
